% I. E. Irodov, Problems in General Physics, user-provided Arihant PDF.
% Problems 241--250; original printed pages 38--39; PDF pages 47--48.
% First draft of Chinese translation with independently worked equations.

\exproblem{241}{38}
一块原本均匀、半径为 $R=20\,\mathrm{cm}$ 的圆盘，
按原书图1.54挖去一个直径为 $R$ 的圆孔：
孔与原盘外缘内切，孔边缘经过原圆心 $O$。
剩余阴影部分质量为 $m=7.3\,\mathrm{kg}$。
求剩余部分对通过其\emph{质心且垂直盘面}的轴的转动惯量。
\exsolution
设原盘面密度为 $\sigma$，孔半径 $R/2$，
孔心在 $x=-R/2$。原盘与被挖去圆片的面积比为 $4:1$，
故其质量分别为 $M=4m/3$、$m_h=m/3$。
取原圆心为 $x=0$，剩余部分质心位置满足
\[
mx_C=M(0)-m_h(-R/2)
\quad\Longrightarrow\quad
\boxed{x_C=R/6}.
\]
对原圆心 $O$ 的惯量可用“完整盘减去孔”计算。
平行轴定理给出被挖圆片对 $O$ 的惯量为
\[
I_{h,O}=\tfrac12m_h(R/2)^2+m_h(R/2)^2
=\tfrac38m_h R^2.
\]
因此
\[
I_O=\tfrac12MR^2-I_{h,O}
=\left(\frac23-\frac18\right)mR^2
=\frac{13}{24}mR^2.
\]
再向剩余部分的质心轴平移：
\[
\boxed{I_C=I_O-mx_C^2
=\left(\frac{13}{24}-\frac1{36}\right)mR^2
=\frac{37}{72}mR^2
\simeq0.150\,\mathrm{kg\,m^2}.}
\]
\excheck 这里的 $m$ 是\emph{挖孔后}剩余部分质量，
不是完整圆盘质量。所得数值与底本书末答案相符。

\exproblem{242}{38}
利用均匀实心球的转动惯量公式，
求质量为 $m$、半径为 $R$ 的均匀薄球壳
关于过球心直径轴的转动惯量。
\exsolution
均匀实心球的质量、转动惯量分别为
$M(R)=\frac43\pi\rho R^3$、
$I(R)=\frac25M(R)R^2=\frac8{15}\pi\rho R^5$。
设把半径从 $R$ 增加到 $R+\mathrm dR$，
新增薄层质量及转动惯量为
\[
\mathrm dm=\frac{\mathrm dM}{\mathrm dR}\mathrm dR
=4\pi\rho R^2\,\mathrm dR,\qquad
\mathrm dI=\frac{\mathrm dI}{\mathrm dR}\mathrm dR
=\frac83\pi\rho R^4\,\mathrm dR.
\]
二者比值
$\mathrm dI/\mathrm dm=\frac23R^2$，
令薄层总质量为 $m$，得到
\[
\boxed{I_{\mathrm{shell}}=\frac23mR^2}.
\]
\excheck 球壳的 $2/3$ 不等于均匀实心球的 $2/5$；
质量都集中于外半径附近时惯量更大。

\exproblem{243}{38}
质量为 $M$、半径为 $R$ 的均匀实心圆柱
可绕固定水平轴无摩擦转动。
圆柱上缠绕轻绳，绳自由端悬挂质量 $m$ 的重物（图1.55）。
系统从静止释放，求：(a) 圆柱角速度 $\omega(t)$；
(b) 包括重物在内的系统总动能 $K(t)$。
\exsolution
以重物向下为正，绳无滑动，故
$a=R\beta$。记绳张力为 $T$，圆柱惯量
$I=\frac12MR^2$。牛顿方程和转动方程为
\[
mg-T=ma,\qquad TR=I\beta=\frac12MRa.
\]
于是 $T=\frac12Ma$，消元得
\[
a=\frac{mg}{m+M/2}=\frac{g}{1+M/(2m)}.
\]
由初始静止和 $\beta=a/R$，
\[
\boxed{\omega(t)=\frac{gt}{R[1+M/(2m)]}}.
\]
重物速度 $v=at$，所以总动能
\[
K=\frac12mv^2+\frac12I(v/R)^2
=\frac12(m+M/2)a^2t^2.
\]
化简后
\[
\boxed{K(t)=\frac{mg^2t^2}{2[1+M/(2m)]}}.
\]
\excheck 两式与原书简答相符；
$m\to\infty$（相对圆柱）时 $a\to g$，
$M\to\infty$ 时 $a\to0$。

\exproblem{244}{38}
麦克斯韦轮的轮轴半径为 $r$，
轮和轮轴总质量为 $m$，绕对称轴转动惯量为 $I$。
两根轻绳对称缠绕于轮轴，
其上端固定在可竖直升降的横梁上。
横梁以恰当的竖直加速度向上运动，使轮的质心高度不变，
同时绳逐渐放开。求：(a) 每根绳的张力；
(b) 横梁向上的加速度。
\exsolution
轮质心始终不作竖直加速，故竖直力平衡
\[
2T-mg=0,\qquad
\boxed{T=\frac{mg}{2}}.
\]
两根绳对轮轴施加同向力矩
$\tau=2Tr= mgr$。
轮绕自身轴的角加速度
$\beta=\tau/I=mgr/I$。
绳与轮轴间不滑动，轮心静止时横梁向上的
线加速度大小等于轮轴表面放绳切向加速度：
\[
\boxed{a_{\rm beam}=r\beta=\frac{mgr^2}{I}}.
\]
\excheck 横梁提供能量使轮加速自转，而轮心高度不变。
该结论的 $I$ 必须包括整个轮与轮轴。

\exproblem{245}{38}
一根质量为 $m$、长度为 $\ell$ 的均匀水平细杆，
可绕通过一端 $A$ 的固定竖轴无摩擦转动。
原来静止，另一端 $B$ 从此时起受到水平恒力 $F$，
力的方向\emph{始终垂直于杆的原始方向}，而非始终垂直于杆。
求角速度随杆转过的角度 $\varphi$ 的关系。
\exsolution
杆关于端点轴的转动惯量
$I_A=m\ell^2/3$。由于力矢量在实验室坐标中
方向固定，杆端沿该方向的位移为
$\Delta y_B=\ell\sin\varphi$，因此外力做功
$W=F\ell\sin\varphi$。
若杆从 $\varphi=0$ 静止开始且没有损耗，
动能定理给出
\[
\tfrac12I_A\omega^2=F\ell\sin\varphi.
\]
故
\[
\boxed{\omega(\varphi)=
\sqrt{\frac{6F\sin\varphi}{m\ell}}}
\]
在从原方向开始加速的 $0\le\varphi\le\pi$
运动段取正向角速度。
\excheck 若错误认为力始终垂直当前杆，
会得到 $\omega^2\propto\varphi$，
与题面“perpendicular to the original position”不符。
底本 PDF 文本抽取可能漏识别答案中的平方根。

\exproblem{246}{39}
如原书图1.56，两个悬挂物体质量分别为
$m_1,m_2$，连接轻绳跨过半径 $R$、
质量 $m$ 的均匀实心圆柱滑轮。
绳与滑轮间无相对滑动，轴无摩擦。
求：(a) 滑轮角加速度；
(b) 两侧绳张力比 $T_1/T_2$。
\exsolution
先取 $m_2>m_1$、物体2向下加速度为 $a$，
绳长约束使物体1向上同加速度。
记滑轮惯量 $I=\frac12mR^2$。列方程
\[
T_1-m_1g=m_1a,\qquad
m_2g-T_2=m_2a,\qquad
(T_2-T_1)R=I\frac aR.
\]
代入张力并化简：
\[
a=\frac{(m_2-m_1)g}{m_1+m_2+m/2},
\quad
\boxed{|\beta|=\frac{|m_2-m_1|g}
{R(m_1+m_2+m/2)}}.
\]
此外
\[
T_1=m_1(g+a)
=\frac{m_1(m+4m_2)g}{2(m_1+m_2+m/2)},
\quad
T_2=m_2(g-a)
=\frac{m_2(m+4m_1)g}{2(m_1+m_2+m/2)}.
\]
故
\[
\boxed{\frac{T_1}{T_2}
=\frac{m_1(m+4m_2)}{m_2(m+4m_1)}}.
\]
\excheck 当 $m_1=m_2$ 时 $\beta=0$ 且张力比为1；
当滑轮质量趋零时亦恢复轻滑轮的相等张力。

\exproblem{247}{39}
质量为 $m_1$ 的物体在摩擦系数为 $k$ 的
水平面上，轻绳跨过质量为 $m$、
可视为均匀圆盘的有质量滑轮，另一端悬挂 $m_2$（图1.57）。
系统从静止释放，且已知 $m_2$ 开始下降。
忽略绳质量与轴摩擦，求最初 $t$ 秒摩擦力对 $m_1$ 做的功。
\exsolution
设 $m_2$ 向下、$m_1$ 向右的加速度均为 $a$，
滑轮半径为 $R$，惯量 $I=\frac12mR^2$。
将三部分的动力学方程相加，绳张力在整体中抵消：
\[
m_2g-km_1g=(m_1+m_2+I/R^2)a.
\]
故
\[
a=\frac{(m_2-km_1)g}{m_1+m_2+m/2}>0.
\]
前 $t$ 秒物体1位移 $s=\frac12at^2$，
滑动摩擦力大小 $km_1g$ 且反向，
\[
W_f=-km_1gs
=\boxed{-\frac{k m_1(m_2-km_1)g^2t^2}
{m+2(m_1+m_2)}}.
\]
\excheck 题面声明 $m_2$ 开始下降，等价于
$ m_2>km_1$ 的滑动模型阈值，故最终功为负。

\exproblem{248}{39}
半径 $R$ 的均匀实心圆柱绕自身轴以初角速度
$\omega_0$ 旋转，随后放入两互相垂直的粗糙墙形成的角落，
如原图1.58。
两接触面摩擦系数同为 $k$，圆柱中心始终保持静止。
求圆柱停止前转过多少圈 $n$。
\exsolution
取原图所示旋向，底部摩擦力水平向左，
左壁摩擦力竖直向上。记底部支持力 $N_b$、
左壁支持力 $N_w$。
质心无平动，水平与竖直方向的力平衡给出
\[
N_w=kN_b,\qquad N_b+kN_w=mg.
\]
所以
$N_b=mg/(1+k^2)$、
$N_w=kmg/(1+k^2)$。
两处滑动摩擦对中心产生同向阻转力矩：
\[
\tau_f=kR(N_b+N_w)
=\frac{k(1+k)mgR}{1+k^2}.
\]
实心圆柱惯量 $I=\frac12mR^2$，
角减速度大小为
$\beta_f=\tau_f/I=
2k(1+k)g/[(1+k^2)R]$。
用 $\omega_0^2=2\beta_f\Delta\varphi$，
并以 $n=\Delta\varphi/(2\pi)$ 换算圈数：
\[
\boxed{n=\frac{(1+k^2)\omega_0^2R}
{8\pi k(1+k)g}}.
\]
\excheck 此结果依赖图示旋向；反向旋转时两面摩擦
竖直、水平分量的组合改变，应重新列平衡方程，
不能不检查受力方向就套用此式。

\exproblem{249}{39}
半径为 $R$ 的均匀圆盘绕竖直中心轴以角速度
$\omega_0$ 旋转，然后轻放在水平粗糙面上。
摩擦系数为 $k$，盘对地的压强近似均匀。
求圆盘继续旋转多久才停止。
\exsolution
圆盘总支持力 $mg$，均匀压强
$p=mg/(\pi R^2)$。半径为 $r$、
宽度为 $\mathrm dr$ 的圆环与地面的接触面积为
$2\pi r\,\mathrm dr$，其摩擦力为
$k p\,2\pi r\,\mathrm dr$，对中心的阻转力矩微元
\[
\mathrm d\tau=r\cdot k p\,2\pi r\,\mathrm dr
=2\pi kp r^2\,\mathrm dr.
\]
积分得 $\tau=\frac23kmgR$。
圆盘惯量 $I=\frac12mR^2$，
角减速度恒为
$\beta=\tau/I=4kg/(3R)$。
于是从角速度 $\omega_0$ 减到零所需时间
\[
\boxed{t_{\rm stop}=\frac{\omega_0}{\beta}
=\frac{3\omega_0R}{4kg}}.
\]
\excheck 若误把所有接触摩擦集中作用在边缘 $r=R$，
会把阻转力矩错误估成 $kmgR$，结果偏小。

\exproblem{250}{39}
一飞轮初始角速度为 $\omega_0$，
受到大小与瞬时角速度平方根成正比的阻转力矩：
$|\tau|=c\sqrt{\omega}$，其中 $c>0$。
求整个减速停止阶段按时间平均的角速度。
\exsolution
设飞轮转动惯量为 $I$，
转动方程
$I\dot\omega=-c\sqrt{\omega}$。
分离变量积分：
\[
2\sqrt{\omega(t)}=
2\sqrt{\omega_0}-\frac cI t.
\]
停止时间为 $T=2I\sqrt{\omega_0}/c$，
因此在 $0\le t\le T$，
\[
\omega(t)=\omega_0\left(1-\frac tT\right)^2.
\]
求时间平均：
\[
\boxed{\bar\omega=\frac1T\int_0^T
\omega_0(1-t/T)^2\,\mathrm dt
=\omega_0\int_0^1(1-u)^2\,\mathrm du
=\frac{\omega_0}{3}}.
\]
\excheck 与角速度衰减成线性时的 $\omega_0/2$ 不同；
阻转系数 $c$ 只决定停转时间，不影响上述无量纲平均值。
